% Не більше 30 % основної частини. Закінчується постановкою задачі.
\section{Передпроєктне обстеження предметної області}

\subsection{Постановка завдання дипломного проєктування}
\todo{...}

\subsection{Аналіз предметної області}
\todo{Визначення, процеси, проблеми.} Приклад рисунка з посиланням: рисунок~\ref{fig:example}.

\begin{figure}
\centering
\includegraphics[width=0.4\textwidth]{example-image}
\caption{Приклад рисунка}
\label{fig:example}
\end{figure}

\subsection{Аналіз існуючих рішень}
\subsubsection{Аналіз відомих програмних продуктів}
\todo{Порівняння аналогів} (таблиця~\ref{tab:analogs}).

\begin{kpitable}{|p{4cm}|p{5cm}|p{5cm}|}{Порівняння аналогів\label{tab:analogs}}{\textbf{Функціонал} & \textbf{Аналог 1} & \textbf{Аналог 2}}
Функція 1 & Так & Ні \\ \hline
\end{kpitable}

\subsubsection{Аналіз відомих алгоритмічних та технічних рішень}
\todo{...} Формула вбудовується в речення після двокрапки:
\begin{equation}
M = \frac{F}{g},
\label{eq:mass}
\end{equation}
\begin{where}
M & маса, кг; \\
F & сила, Н; \\
g & прискорення вільного падіння, м/с$^2$.
\end{where}

Нумерують лише формули, на які далі є посилання, як-от у формулі~(\ref{eq:mass}).

\subsection{Аналіз та моделювання бізнес-процесів}
\todo{...}

\subsection{Висновки до розділу}
\todo{...}
